% Nouveau EGA V (cote 25): archive cote 25, site page 116 (PDF page 117).
% typescript. Draft compared: none.
% AI-transcribed by Claude Opus 5.5 (claude-opus-5-5, adaptive thinking, effort high), Message Batches API, batch ega25b (msgbatch_01A7n5YuXubUPNK1zsepi8RG), 2026-11-13; checked against both supplied scan parts; see CHECK_DECISIONS.tsv.
\scanpage{25/116}
\mspage{19}

points fixes, elle est déterminée \var{\xout{uniquement}}{}{typewriter strike-out}{ega25_25_116-v0001} par le morphisme f:X $\hand{\longrightarrow}$ P
associé. Noter d'ailleurs qu'on \var{\deleted{\var{\xout{\illegible{1}}}{}{typewriter strike-out}{ega25_25_116-v0003}}}{}{deletion}{ega25_25_116-v0002}obtient évidemment ainsi une correspon
dance biunivoque entre systêmes linéaires de diviseurs de X/S paramétré
par $\hand{\check{P}}$ , et morphismes f:X $\hand{\longrightarrow}$ P tel que \var{\deleted{\var{\xout{\illegible{1}}}{}{typewriter strike-out}{ega25_25_116-v0005}}}{}{deletion}{ega25_25_116-v0004}$(fx_S\mathrm{id})^{-1}(\underline{H})$ soit un
diviseur\var{\deleted{ de Cartier}}{}{deletion}{ega25_25_116-v0006} \var{\xout{sur}}{}{typewriter strike-out}{ega25_25_116-v0007} relatif sur $Xx_S\overset{\hand{\vee}}{P}$ par rapport à $\overset{\hand{\vee}}{P}$ . Ce\var{\over{\unsure{ci}}{tte}}{tte}{writing over}{ega25_25_116-v0008} \var{\deleted{se}}{}{deletion}{ega25_25_116-v0009} \var{\inserted{condition}}{condition}{insertion}{ega25_25_116-v0010}
\var{\inserted{se}}{se}{insertion}{ega25_25_116-v0011} vérifie aussi fibre par fibre, et on obtient :

\uline{Propositio\unsure{n|h}} \var{\over{\unsure{1|2}}{3}}{3}{writing over}{ega25_25_116-v0012}.\hand{1}.On a une correspondance biunivoque entre les \var{\xout{morphismes}}{}{typewriter strike-out}{ega25_25_116-v0013}
systêmes linéaires \var{\inserted[typed]{sans points fixes}}{sans points fixes}{insertion}{ega25_25_116-v0014} de diviseurs sur X/S paramétrés par $\hand{\check{P}}$ , et les morphis
mes f: X $\hand{\longrightarrow}$ P ayant la propriété suivante: pour tout $s\hand{\in} S$, d\var{\over{e}{é}}{é}{writing over}{ega25_25_116-v0015}signant
par $\bar{k}$ une clôture algébrique de k(s), et pour tout cycle premier associé
X' de $X_{\bar{k}}$ , $f(X')\hand{\subset} P_{\bar{k}}$ n'est contenu dans aucun hyperplan de $P_{\bar{k}}$ .
\var{\xout{(NB}}{}{typewriter strike-out}{ega25_25_116-v0016}\hand{(}Lorsque X est à fibres géométriquement \var{\xout{\unsure{intégres|intègres}}}{}{typewriter strike-out}{ega25_25_116-v0017} \var{\inserted[typed]{intègres}}{intègres}{insertion}{ega25_25_116-v0018}, ceci \var{\xout{se simplifie}}{}{typewriter strike-out}{ega25_25_116-v0019}
\var{\xout{et}}{}{typewriter strike-out}{ega25_25_116-v0020} s'énoncé simplement en disant que $f(X_{\bar{k}})$ n'est contenu dans aucun
hyperplan de $P_{\bar{k}}$ \hand{)}.

En général (i.e. \var{\xout{en}}{}{typewriter strike-out}{ega25_25_116-v0021} lorsque Z $\hand{\neq}\hand{\emptyset}$ ) on ne peut plus affirmer que
la connaissance de f détermine \unsure{D|l}a famille de diviseurs. Le cas le plus
trivial est celui où $\hand{\check{P}}$ = S \var{\inserted[typed]{est de dim relative 0}}{est de dim relative 0}{insertion}{ega25_25_116-v0022}, alors la donnée d'un systême linéaire de\hand{.}
diviseurs de X/S paramétrée par S équivaut à celle d'un diviseur de
Cartier relatif D sur X par rapport à S; \var{\xout{Tout l'app\illegible{1}}}{}{typewriter strike-out}{ega25_25_116-v0023} le morphisme
associé est la projection X-D $\hand{\longrightarrow}$ S , et on voit \var{\xout{(même en}}{}{typewriter strike-out}{ega25_25_116-v0024} que la connais
sance de ce morphisme (y compris celle de son domaine de définition)
ne détermine pas D. Dans ce cas également, supposant pour simplifier X
réduit, le domaine de définition de \var{\xout{ce morphisme}}{}{typewriter strike-out}{ega25_25_116-v0025} f \var{\xout{n'est pas}}{}{typewriter strike-out}{ega25_25_116-v0026} considéré
comme application rationelle de X dans P = S n'est pas X-D , mais X.
Pour éliminer ce genre de phénomènes déplaisants, il y a lieu de se bor-%
ner aux systêmes \unsure{h|l}i\unsure{\'n|n}éaires de diviseurs ``\uline{sans composante fixe}'' . De
façon générale, si S=Spec(k), étaht donné une famille (non néc linéaire)

\var{\margin{116}}{\placedmargin{116}}{marginal note placed}{ega25_25_116-v0027}
